% Requires in the paper preamble:
% \usepackage{tikz}
% \usetikzlibrary{matrix,positioning}
% \definecolor{ObservedCell}{HTML}{DCEAD7}
% \definecolor{HeldOutCell}{HTML}{F4D0CC}

\begin{figure}[t]
    \centering
    \begin{tikzpicture}[
        cell/.style={
            draw=black!35,
            minimum height=8.5mm,
            minimum width=25mm,
            align=center,
            font=\footnotesize
        }
    ]
        \matrix (m) [
            matrix of nodes,
            nodes in empty cells,
            column sep=-\pgflinewidth,
            row sep=-\pgflinewidth,
            row 1/.style={nodes={font=\footnotesize\bfseries,draw=none}},
            column 1/.style={nodes={font=\footnotesize\bfseries,draw=none}}
        ] {
            & Other training objects & Green block \\
            \emph{smallest} & |[cell,fill=ObservedCell]| Observed & |[cell,fill=ObservedCell]| Observed \\
            \emph{largest}  & |[cell,fill=ObservedCell]| Observed & |[cell,fill=HeldOutCell,dashed,draw=black!70]| \textbf{Withheld} \\
        };

        \node[
            below=4mm of m,
            align=center,
            font=\scriptsize,
            text width=0.86\columnwidth
        ] (masked) {The green block and the corresponding pick-and-place trajectory remain in training under a direct-name prompt, e.g., ``Move the green block into the bowl.''};

        \draw[->,semithick,black!60]
            (m-3-3.south) -- ++(0,-2mm) -| (masked.north east);
    \end{tikzpicture}
    \caption{Schematic of one missing cell in Size Recognition. Training explicitly pairs \emph{largest} with other objects and pairs the green block with \emph{smallest}, but never pairs \emph{largest} with the green block. The withheld concept--argument binding is evaluated even though the concept, object identity, and physical trajectory each remain represented in training.}
    \label{fig:missing_cell_matrix}
\end{figure}

